\chapter{Software testing}

Tests define the expected behavior of code, and detect when the code doesn't match that expected behavior.

One useful analogy for software testing comes from the biosciences.
Think for a moment about the rapid COVID-19 tests that we all came to know during the pandemic.
These tests had two lines, one of which was a \textit{control} line; if this line didn't show up, then that meant that the test was not functioning as expected.
This is known as a \textit{positive control} because it assesses the test's ability to identify a positive response\footnote{This is slightly inaccurate, because a true positive control would contain the actual virus.
It would be more precise to call it a ``procedural control'' but these seem to be also referred to as ``positive controls'' so I am sticking with the more understandable terminology here.}.
Other tests also include \textit{negative controls}, which ensure that the test returns a negative result when it should.

By analogy, we can think of software tests as being either positive or negative controls for the expected outcome of a software component.
A positive test assesses whether, given a particular valid input, the component returns the correct output.
A negative test assesses whether, in the absence of valid input, the component correctly returns the appropriate error message or null result.

\section{Why use software tests?}

The most obvious reason to write tests for code is to make sure that the answers that the code gives you are correct.
This becomes increasingly important as AI assistants write more of the code, to the degree that testing is becoming \textit{at least as important} as code generation as a skill for generating good scientific code.
But creating correct code is far from the only reason for writing tests.

A second reason for testing was highlighted in our earlier discussion of test-driven development.\index{test-driven development}
Tests can provide the coder with a measure of task completion; when the tests pass, the job is done, other than refactoring the code to make it cleaner and more robust.\index{refactoring}
Writing tests make one think harder about what exactly they want/need the code to do, and to specify those goals in as clear a way as possible.
Focusing on tests can help keep the coder's ``eyes on the MVP prize'' and prevent generating too much extraneous code (``gold plating'').

A third reason to write tests is that they can help drive modularity in the code.\index{modularity}
It's much easier to write tests for a simple function that does a single thing than for a complex function with many different roles.
Testing can also help drive modularity by causing one to think more clearly about what a function does when developing the test; the inability to easily write a test for a function can suggest that the function might be overly complex and should be refactored.
In this way, writing tests can give us useful insights into the structure of the code.

A final reason to write tests is that they make it much easier to make changes to the code.
Without a robust test suite, one is always left worried that changing some aspect of the code will have unexpected effects on its former behavior (known as a ``regression'').
Tests can provide you with the comfort you need to make changes, knowing that you will detect any untoward effects your changes might have.
This includes refactoring, where the changes are not meant to modify the functionality but simply to make the code more robust and readable.

\section{Types of tests}

There are many different types of tests that are discussed in software engineering, but here I will focus on two types that are most important for scientific code: unit tests and regression tests.\index{testing!unit}

\subsection{Unit tests}

Unit tests are the bread and butter of software testing.
They are meant to assess whether individual software components (in the case of Python, functions, classes, and methods) perform as expected.\index{Python}
This includes both assessing whether the component performs as it is supposed to perform given a particular input, but also assessing whether it performs correctly under boundary conditions or problematic conditions, where the correct response is often to raise an exception.
A major goal of unit testing in the latter case is preventing ``garbage in, garbage out'' behavior.
For example, say that we are testing a function that takes in two matrices, and that the size of these matrices along their first dimension is assumed to match.
In this case, we would want to test to make sure that if the function is provided with two matrices that mismatch in their first dimension, the function will respond by raising an exception rather than by giving back an answer that is incorrect or nonsensical (such as \textit{NaN}, or ``not a number'').
That is, we want to aim for ``garbage in, exception out'' behavior.

\subsection{Integration tests}\index{testing!integration}

As the name suggests, an integration test assesses whether the entire application works as it should, integrating all of the components that were tested via unit testing.

One simple type of integration test is a ``smoke test''.\index{testing!smoke test}
This name apparently derives from the computer hardware industry, where one often performs an initial sanity test on an electronic component by plugging it in and seeing if it smokes.
In coding, a smoke test is a simple sanity check meant to ensure that the entire application runs without crashing.
This is usually accomplished by running a top-level function that exercises the entire application.
Smoke tests are useful for quickly identifying major problems, but they don't actually test whether the application performs its function correctly.
They can be especially useful for large applications, where the full test suite may take hours to run.
An initial smoke test can determine whether something is broken downstream, saving lots of wasted testing time.

Full integration tests assess the function of the entire application; one can think of them as unit tests where the unit is the entire application.
Just as with unit tests, we want integration tests that both confirm proper operation under intended conditions, as well as confirming proper behavior (such as exiting with an error message) under improper conditions.

\section{The anatomy of a test}

A test is generally structured as a function that executes without raising an exception as long as the code behaves in an expected way.
Let's say that we want to generate a function that returns the escape velocity of a planet:

\lstinputlisting[style=Python]{../src/bettercode/testing/escape_velocity_v1.py}

We can then generate a test to determine whether the value returned by our function matches the known value for a given planet:

\lstinputlisting[style=Python, firstline=8, lastline=16]{../tests/testing_chapter/test_escape_velocity.py}

\subsection{Assertions}

The statement at the end of the function is the critical component known as an \textit{assertion}, so let's break it down to understand it.\index{assertion} The \texttt{assert} command follows the pattern: 

\begin{lstlisting}[style=Pythonshort]
assert <statement returning a truth value>, <message>
\end{lstlisting}

If the statement following \texttt{assert} returns \texttt{True}, then nothing happens.  If the statement returns \texttt{False}, then the \texttt{assert} command raises an \texttt{AssertionError} exception with the specified message; if the statement itself raises an Exception, then that exception is passed forward by the command.  This is an essential tool for testing: We assert all of the things that should be true, such that if there are no exceptions raised then that means the tests have passed.

We can run this using \textit{pytest} (more about this later), which tells us that the test passes:\index{pytest}

\lstinputlisting[style=shell]{generated/testing/pytest_escape_velocity_pass.txt}

If the returned value didn't match the known value (within a given level of tolerance, which is handled by \texttt{np.allclose()}), then the assertion will fail and raise an exception, causing the test to fail.\index{assertion}. 
For example, if we had mis-specified the expected value as 1186.0, we would have seen an error like this:

\lstinputlisting[style=shell]{generated/testing/pytest_escape_velocity_typo_fail.txt}

When testing for equality of numbers (except for integers), it is essential to use a comparison method that allows for a degree of tolerance in the comparison.\index{tolerance}  For example, it might seem like the following assertion should pass, but it doesn't:

\lstinputlisting[style=repl, numbers=none]{generated/testing/repl_float_assert_fail.txt}

This happens because the result of \textit{0.1 + 0.2} is very close, but not exactly the same as \textit{0.3} due to the nature of how floating point numbers are represented:

\lstinputlisting[style=repl, numbers=none]{generated/testing/repl_float_value.txt}

Instead, we want to compare the two numbers using a method that allows for a degree of \textit{tolerance} or slight differences between the numbers at the smallest decimal places.  A commonly used tool is the \texttt{np.allclose} function:

\lstinputlisting[style=repl, numbers=none]{generated/testing/repl_float_allclose.txt}

Used this way the function uses default tolerances.  There are two kinds of tolerance: \textit{Absolute tolerance} is simply the absolute difference between the values (default: 1e-08), while \textit{Relative tolerance} is proportional to the size of one of the items (the second one for texttt{np.allclose}) (default: 1e-05).  This function uses a comparison that combines the relative and absolute tolerances.  What's important to keep in mind is that tolerance relates to the size of the numbers, so when working with either very small or very large numbers it's often necessary to adjust the tolerance to prevent cases where the test fails due to floating point errors or passes even when the value is incorrect.

\subsection{Checking for exceptions}

It's also important to make sure that an exception is raised whenever it should be.
For example, the version of the \texttt{escape\_velocity()} function above did not check to make sure that the mass and radius arguments had positive values, which means that it would give nonsensical results when passed a negative mass or radius value.
To address this we should add code to the function that causes it to raise an exception if either of the arguments is negative:

\lstinputlisting[style=Python]{../src/bettercode/testing/escape_velocity_v2.py}

We can then specify a test that checks whether the function properly raises an exception when passed a negative value.
To do this we can use a feature of the \textit{pytest} package (\texttt{pytest.raises}) that passes only if the specified exception is raised:

\lstinputlisting[style=Python, firstline=19, lastline=24]{../tests/testing_chapter/test_escape_velocity.py}

\section{When to write tests}

Too often researchers decide to write tests after they have written an entire codebase.
Having any tests is certainly better than having no tests, but integrating testing into one's development workflow from the start can help improve the development experience and ultimately lead to better and more maintainable software.
In Chapter 3 I mentioned the idea of \textit{test-driven development}, which I outline in more detail below, but I first discuss a simple approach to introducing testing into the development process.

\subsection{Bug-driven testing: Any time you encounter a bug, write a test}

An easy way to introduce testing into the development process is to write a new test any time one encounters a bug, which I refer to as \textit{bug-driven testing}.
This makes it easy to then work on fixing the bug, since the test will determine when the bug has been fixed.
In addition, the test will detect if future changes reintroduce the bug.

As an example, take the following function:

\lstinputlisting[style=Python]{../src/bettercode/testing/find_outliers_v1_buggy.py}

This code works to properly identify outliers:

\lstinputlisting[style=repl, numbers=none]{generated/testing/repl_find_outliers_normal.txt}

However, it fails due to a division by zero if all of the values are equal:

\lstinputlisting[style=repl]{generated/testing/repl_find_outliers_zero_div.txt}

Our intended behavior if all of the values are equal is to return an empty list, since there are by definition no outliers.
But before we do this, let's create a couple of tests to check for the intended behavior and provide useful error messages if the test fails:

\lstinputlisting[style=Python]{../tests/testing_chapter/test_find_outliers.py}

Running this with the original function definition, we see that it fails:

\lstinputlisting[style=shell]{generated/testing/pytest_find_outliers_bug_fail.txt}

We can now fix the code by returning an empty list if zero standard deviation is detected:

\lstinputlisting[style=Python, firstline=25, lastline=28]{../src/bettercode/testing/find_outliers_v2_fixed.py}

Here we add a comment to explain the intention of the statement.
Running the tests now will show that the problem is fixed:

\lstinputlisting[style=shell]{generated/testing/pytest_find_outliers_fixed_pass.txt}

Now we can continue coding with confidence that if we happen to accidentally reintroduce the bug, it will be caught.

\section{The structure of a good test}

A commonly used scheme for writing a test is ``given/when/then'':

\begin{itemize}
\item given some particular situation as background
\item when something happens (such as a particular input)
\item then something else should happen (such as a particular output or exception)
\end{itemize}

Importantly, a test should only test one thing at a time.
This doesn't mean that the test should necessarily only test for one specific error at a time; rather, it means that the test should assess a specific situation (``given/when''), and then assess all of the possible outcomes that are necessary to ensure that the component functions properly (``then'').
You can see this in the test for zero standard deviation that we generated in the earlier example, which actually tested for two conditions (the intended value being present in the list, and the list having a length of one) that together define the condition that we are interested in testing for.

How do we test that the output of a function is correct given the input? There are different answers for different situations:

\begin{itemize}
\item \textit{commonly known answer}: Sometimes we possess inputs where the output is known.  For example, if we were creating a function that computes the circumference of a circle, then we know that the output for an input radius of 1 should be 2 * pi.  This is generally only the case for very simple functions.
\item \textit{reference implementation}: In other cases we may have a standard implementation of an algorithm that we can compare against.  While in general it's not a good idea to reimplement code that already exists in a standard library, in come cases we may want to extend existing code but also check that the basic version still works as planned.
\item \textit{parallel implementation}: Some times we don't have a reference implementation, but we can code up another parallel implementation to compare our code to.  It's important that this isn't just a copy of the code used in the function; in that case, it's really not a test at all!
\item \textit{behavioral test}: Sometimes the best we can do is to run the code repeatedly and ensure that it behaves as expected on average.  For example, if a function outputs a numerical value and we know the expected distribution of that value given a particular input, we can ensure that the result matches that distribution with a high probability.  Such \textit{probabilistic tests} are not optimal in the sense that they can occasionally fail even when the code is correct, but they are sometimes the best we can do.
\end{itemize}

\subsection{Test against the interface, not the implementation}

A good test shouldn't know about the internal implementation details of the function that it is testing, and changes in the internal code that do not modify the input-output relationship should not affect the test.
That is, from the standpoint of the test, a function should be a ``black box''.\index{black (formatter)}

The most common way in which a test can violate this principle is by accessing the internal variables of a class that it is testing.
For example, we might generate a class that performs a scaling operation on a numpy matrix:

\lstinputlisting[style=Python]{../src/bettercode/testing/simple_scaler.py}

We could write a test that checks the values returned by the \texttt{fit\_transform()} method, treating the the class as a black box:

\lstinputlisting[style=Python, firstline=6, lastline=13]{../tests/testing_chapter/test_simple_scaler.py}

Alternatively one might use knowledge of the internals of the class to test the transformed value:

\lstinputlisting[style=Python, firstline=16, lastline=23]{../tests/testing_chapter/test_simple_scaler.py}

Both of these tests pass against the class definition shown above.
However, if we were to change the way that the transformation is performed (for example, we decide to use the \texttt{StandardScaler} function from \texttt{scikit-learn} instead of writing our own), then the implementation-aware tests are likely to fail unless the same internal variable names are used.\index{scikit-learn}
In general we should only interact with a function or class via its explicit interfaces.

\subsection{Tests should be independent}

In scientific computing it's common to compose many different operations into a workflow.
If we want to test the workflow, then the tests of later steps in the workflow must necessarily rely upon earlier steps.
We could in theory write a set of tests that operate on a shared object, but the tests would fail if executed in an incorrect order, even if the code was correct.
Similarly, a failure on an early test would cause cascading failures in later tests, even if their code was correct.
The use of ordered tests also prevents the parallel execution of tests, which may slow down testing for complex projects.
For these reasons, we should always aim to create tests that can be executed independently.

Here is an example where coupling between tests could cause failures.\index{coupling}
First we generate two functions that make changes in place to a data frame:

\lstinputlisting[style=Python]{../src/bettercode/testing/dataframe_ops.py}

In this case, the \texttt{get\_initials()} function relies upon the \texttt{split\_names()} function having been run, since otherwise the necessary columns won't exist in the data frame.
We can then create tests for each of these, and a data frame that they can both use:

\lstinputlisting[style=Python, firstline=5]{../tests/testing_chapter/_demos/test_dataframe_ops_coupled.py}

These tests run correctly, but the same tests fail if we change their order such that \texttt{test\_get\_intials()} runs first, because the necessary columns (\texttt{firstname} and \texttt{lastname}) have not yet been created.

One simple way to deal with this is to set up all of the necessary structure locally within each test:

\lstinputlisting[style=Python, firstline=5]{../tests/testing_chapter/test_dataframe_ops.py}

For simple functions like these this would not cause too much computational overhead, but for computationally intensive functions we would like to be able to reuse the results from the first time each function is run.
In a later section we will discuss the use of \textit{fixtures} which allow this kind of reuse across tests while avoiding the ordering problems that we saw above when using a global variable across tests.\index{global variable}

\section{Testing frameworks}

One could write tests without the help of any specialized packages, but we generally use a testing framework to help automate the activities of testing.
There are several testing frameworks for Python; we will focus on the popular and flexible \textit{pytest} framework.

We will start with a very simple example: a function that generates the Euclidean distance between two points.
Our coding assistant generates the following for us based on the prompt in the comment:

\lstinputlisting[style=Python]{../src/bettercode/testing/distance.py}

Now we would like to generate some tests for this code to make sure that it works properly.
If we ask Copilot to generate some tests, it does a seemingly decent job:

\lstinputlisting[style=Python]{../tests/testing_chapter/_demos/test_distance_v1_failing.py}

Now that we have our tests, we can run them using the \texttt{pytest} command:

\begin{lstlisting}[style=shell, numbers=none]
$ pytest tests/test_distance.py
\end{lstlisting}

This command will cause pytest to search (by default) for any files named \texttt{test\_*.py} or \texttt{*\_test.py} in the relevant path, and the select any functions whose name starts with the prefix ``test''.
Running those tests, we get an error:

\lstinputlisting[style=shell]{generated/testing/pytest_distance_fail.txt}

Here we see that the value returned by our function is different from the one expected by the test; in this case, the test value generated by Copilot is incorrect.
In our research, it was not uncommon for AI coding tools to generate correct tests structures but use incorrect test values. 
Test values should always be checked by a domain expert.
Once we fix the expected value for that test (the square root of 89), then we can rerun the tests and see that they have passed:

\lstinputlisting[style=shell]{generated/testing/pytest_distance_pass.txt}

\subsection{Potential problems with AI-generated tests}

If we are going to rely upon AI tools to generate our tests, we need to be sure that the tests are correct.
One of my early forays into AI-driven test generation uncovered an interesting example of how this can go wrong.

In our early project that examined the performance of GPT-4 for coding \citep{Poldrack:2023aa}, one of the analyses that we performed first asked GPT-4 to generate a set of functions related to common problems in several scientific domains, and then to generate tests to make sure that the function performed correctly.\index{GPT-4}
One of the functions that was generated was the escape velocity function shown above, for which GPT-4 generated the following test:

\lstinputlisting[style=Python, firstline=6, lastline=21]{../tests/testing_chapter/_demos/test_escape_velocity_gpt4.py}

When we run this test (renaming it \texttt{test\_escape\_velocity\_gpt4}), we see that one of the tests fails:

\lstinputlisting[style=shell]{generated/testing/pytest_escape_velocity_gpt4_fail.txt}

It seems that the first two assertions pass but the third one, for Jupiter, fails.
This failure took a bit of digging to fully understand.
In this case, the code and test value are both correct, depending on where you stand on Jupiter! The problem is that planets are \textit{oblate}, meaning that they are slightly flattened such that the radius around the equator is higher than at other points.
NASA's Jupiter fact sheet (nssdc.gsfc.nasa.gov/planetary/factsheet/jupiterfact.html) claims an escape velocity of 59.5 km/s, which seems to be the source of the test value.
This is correct when computed using the equatorial radius of 71492 km.
However, the radius given for Jupiter in GPT-4's test (69911 km) is the volumetric mean radius rather than the equatorial radius, and the value generated by the code (60.2 km/s) is correct when computed using the volumetric mean radius.
Thus, the test failed not due to any problems with the code itself, but due to a mismatch in assumptions regarding the combination of test values.
This example highlights the importance of understanding and checking the tests that are generated by AI coding tools based on domain knowledge.

\section{Test-driven development and AI-assisted coding}

Here I will dive into a more realistic example of an application that one might develop using AI assistance, specifically looking at how we could develop the application using a test-driven development (TDD) approach.
I will develop a Python application that takes in a query for the PubMed database and returns a data frame containing the number of database records matching that query for each year.
I start by decomposing the problem and sketching out the main set of functions that we will need to develop, with understandable names for each:

\begin{itemize}
\item \texttt{get\_PubMedIDs\_for\_query}: A function that will search PubMed for a given query and return a list of PubMed IDs
\item \texttt{get\_record\_from\_PubMedID}: A function that will retrieve the record for a given PubMed ID
\item \texttt{parse\_year\_from\_PubMed\_record}: A function that will parse a record to extract the year of publication
\item A function that will summarize the number of records per year
\item The main function that will take in a query and return a data frame with the number of records per year for the query
\end{itemize}

We can then get started by creating \texttt{get\_PubMedIDs\_for\_query}.
We could use the \texttt{Biopython.Entrez} module to perform this search, but Biopython is a relatively large module that could introduce technical debt.\index{Biopython}
Instead, we will directly retrieve the result using the Entrez API and the built-in \texttt{requests} module.
Note that for all of the code shown here I will not include docstrings, but they are available in the code within the repository.

If we are using the TDD approach, we would first want to develop a set of tests to make sure that our function is working correctly.
The following three tests specify several different outcomes that we might expect.
First, we give a query that is known to give a valid result, and test whether it in fact gives such a result:

\lstinputlisting[style=Python, firstline=35, lastline=42]{../tests/testing_chapter/test_textmining.py}

Second, we give a query with a known empty result, and make sure it returns an empty list:

\lstinputlisting[style=Python, firstline=45, lastline=52]{../tests/testing_chapter/test_textmining.py}

With the minimal tests in place, we then move to writing the code for the module.
We first create an empty function to ensure that the tests fail:

\begin{lstlisting}[style=Pythonshort]
def get_PubMedIDs_for_query(query: str, 
                            retmax: int = None,
                            esearch_url: str = None) -> list:
    return None
\end{lstlisting}

The test result shows that all of the tests fail:

\begin{lstlisting}[style=shell]
$ python -m pytest -v test_textmining.py
============================= test session starts =============================
...
test_textmining.py::test_get_PubMedIDs_for_query_check_valid FAILED [ 50%]
test_textmining.py::test_get_PubMedIDs_for_query_check_empty FAILED [100%]

=================================== FAILURES ==================================
___________________ test_get_PubMedIDs_for_query_check_valid ___________________

ids = None

    def test_get_PubMedIDs_for_query_check_valid(ids):
>       assert isinstance(ids, list)
E       assert False
E        +  where False = isinstance(None, list)

test_textmining.py:32: AssertionError
___________________ test_get_PubMedIDs_for_query_check_empty ___________________

    def test_get_PubMedIDs_for_query_check_empty():
        query = "friston -k AND 'fizzbuzz'"
        ids = get_PubMedIDs_for_query(query)
>       assert len(ids) == 0
               ^^^^^^^^
E       TypeError: object of type 'NoneType' has no len()

test_textmining.py:39: TypeError
=========================== short test summary info ==============================
FAILED test_textmining.py::test_get_PubMedIDs_for_query_check_valid - assert False
FAILED test_textmining.py::test_get_PubMedIDs_for_query_check_empty - 
TypeError: object of type 'NoneType' has no len()
============================== 2 failed in 0.12s =================================
\end{lstlisting}

Now we work with Copilot write the code to make the tests pass:

\lstinputlisting[style=Python, firstline=1, lastline=45]{../src/bettercode/testing/textmining.py}

Note that we have split parts of the functionality into separate functions in order to make the code more understandable.
Running the tests, we see that both of them pass.
Assuming that our tests cover all possible outcomes of interest, we can consider our function complete.
We can also add additional tests to cover additional functions that we generated; we won't go into the details here, but you can see them on the project Github repo in the file \texttt{tests/test\_textmining.py}.

\section{Test coverage}

It can be useful to know if there are any portions of our code that are not being exercised by our tests, which is known as \textit{code coverage}.\index{code coverage}
The \textit{pytest-cov} extension for the \textit{pytest} testing package can provide us with a report of test coverage for these tests:

\lstinputlisting[style=shell]{generated/testing/pytest_textmining_coverage.txt}

This report shows that of the 30 statements in our code, one of them is not covered by the tests.
When we look at the missing code (denoted as being on line 70), we see that the missing line is this one from \texttt{get\_idlist\_from\_response}:

\begin{lstlisting}[style=Pythonshort]
else:
        # raise an exception if the search didn't return a usable response
        raise ValueError("Bad request")
\end{lstlisting}

Since none of our test cases caused a bad request to occur, this line never gets executed in the tests.
We can address this by adding a test that makes sure that an exception is raised if an invalid base url is provided.
To check for an exception, we need to use the \texttt{pytest.raises} context manager:\index{context manager}

\lstinputlisting[style=Python, firstline=55, lastline=62]{../tests/testing_chapter/test_textmining.py}

After adding this test, we see that we now have 100\% coverage.
It's important not to get too hung up on test coverage; rather than always aspiring to 100\% coverage, it's important to make sure that the most likely possible situations are tested.
Just because you have 100\% coverage doesn't mean that your code is perfectly tested, since there could always be situations that you haven't checked for.
And spending too much time testing for unlikely problems can divert your efforts from other more useful activities.

\section{Test fixtures}

Sometimes we need to use a the same data for multiple tests.
Rather than duplicating potentially time-consuming processes across each of the tests, it is often preferable to create a single instance of the object that can be used across multiple tests, which is known as a \textit{test fixture}.\index{testing!fixture}
This also helps maintain isolation between tests, since the order of tests shouldn't matter if an appropriate fixture is generated as soon as it's needed.

For our example above, it's likely that we will need to reuse the list of pubmed IDs from the search to perform various tests on the subsequent functions.
We can create a single version of this list of IDs by creating a fixture.
In the \textit{pytest} framework we do this using a special Python operator called a \textit{decorator}, which is denoted by the symbol \texttt{@} as a prefix.\index{decorator}
A decorator is function that takes another function as input, modifies its functionality, and returns another function; you don't need to understand in detail how decorators work for this particular usage.
To refactor our tests above, we would first create the fixture by decorating the function that generates the fixture with the \texttt{@pytest.fixture} decorator, setting the \texttt{scope} variable to ``session'' so that the fixture is only generated once within the session:

\lstinputlisting[style=Python, firstline=14, lastline=18]{../tests/testing_chapter/test_textmining.py}

We can then refactor our tests for a valid query to use the fixture by passing it as an argument to the test function:

\begin{lstlisting}[style=Pythonshort]
def test_get_PubMedIDs_for_query_check_valid(ids):
    assert isinstance(ids, list)
    assert len(ids) > 0
\end{lstlisting}

The result is the same, but we now have a set of ids that we can reuse in subsequent tests, so that we don't have to make repeated queries.
It's important to note that when one is using a session-scoped fixture, any modifications of the fixture by subsequent tests will persist, which will break the isolation between tests.
We could prevent this by removing the \texttt{scope="session"} argument, which would then default to the standard scope which is within a specific function.
If you wish to use session-scoped fixtures and need to modify them within the test function, then it is best to first create a copy of the fixture object (e.g., \texttt{my\_ids = ids.copy()}) so that the global fixture object won't be modified.

\section{Mocking}

Sometimes tests require infrastructure that is outside of the control of the tester.
In the example above, we are assuming that the PubMed API is working correctly for our tests to run; if we were to try to run these tests without an internet connection, they would fail.
In other cases, code may rely upon a database system that may or may not exist on a particular system.
In these cases, we can create a \textit{mock} object that can stand in for and simulate the behavior of the system that the code needs to interact with.

In our example, we want to create a mock response that looks sufficiently like a response from the real API to pass our tests.\index{testing!mocking}
Using \textit{pytest}'s \texttt{monkeypatch} fixture, we can temporarily replace the real requests.get function with our own fake function that returns a predictable, controlled response.
We first need to create a class that can replace the \texttt{requests.get} call in \texttt{get\_PubMedIDs\_for\_query}, replacing it with a mock version that outputs a fixed simulacrum of an API response via its \texttt{.json()} method.

\begin{lstlisting}[style=Python]
class MockPubMedResponse:
    status_code = 200

    def json():
        return {
            'header': {'type': 'esearch', 'version': '0.3'},
            'esearchresult': {
                'count': '2',
                'retmax': '20',
                'retstart': '0',
                'idlist': ['39312494', '39089179']
            }
        }
\end{lstlisting}

We now insert this mock response for the standard \texttt{requests.get} call within the test.
In my initial attempt, I created created a fixture based on the mocked response and then tested that fixture:

\begin{lstlisting}[style=Python]
@pytest.fixture
def ids_mocked(monkeypatch):

    def mock_get(*args, **kwargs):
        return MockPubMedResponse()

    # apply the monkeypatch for requests.get to mock_get
    monkeypatch.setattr(requests, "get", mock_get)

    query = "friston-k AND 'free energy'"
    ids = get_PubMedIDs_for_query(query)
    return ids

def test_get_PubMedIDs_for_query_check_valid_mocked(ids_mocked):
    assert isinstance(ids_mocked, list)
    assert len(ids_mocked) == 2
\end{lstlisting}

Turning off my network connection shows that the mocked test passes, while the tests that require connecting to the actual API fail.
However, my usual code review (using Google's Gemini 2.5 Pro) identified a problem with this fixture: it conflates the setup (creating the mock API) with the execution of the function that uses the mock API.\index{code review}
A better approach is move the function execution out of the fixture and into the test:

\lstinputlisting[style=Python, firstline=65, lastline=100]{../tests/testing_chapter/test_textmining.py}

Note that while mocking can be useful for testing specific components by saving time and increasing robustness, integration tests and smoke tests should usually be run without mocking, in order to catch any errors that arise through interaction with the relevant components that are being mocked.
In fact, it's always a good idea to have tests that specifically assess the usage of the external service and the system's response to failures in that service (e.g., by using features of the testing framework that allow one to shut down access to the network).

\section{Parametrized tests}

Often a function needs to accept a range of inputs that can result in different behavior, and we want to test each of the possible inputs to ensure that the function works correctly across the range.
When the different inputs are known, one way to achieve this is to use a \textit{parameterized test}, in which the test is repeatedly run across combinations of different possible values.

For our example, let's move forward and develop the function \texttt{parse\_year\_from\_PubMed\_record} to extract the year from PubMed records, which can differ in their structure.
We first need to develop the function \texttt{get\_record\_from\_PubMedID} to retrieve a record based on a PubMed ID.
Following our TDD approach, we first develop two simple tests: one to ensure that it returns a non-empty dictionary for a valid PubMed ID, and one to ensure that it raises an exception for an invalid PubMed ID.
We also need to create empty functions so that they can be imported to run the (failing) tests:

\begin{lstlisting}[style=Python]
def get_record_from_PubMedID(pmid: str) -> dict:
    pass

def parse_year_from_PubMed_record(pubmed_record: dict) -> int:
    pass
\end{lstlisting}

Here are the initial tests; note that writing these tests requires a bit of knowledge about the expected structure of a PubMed record.
We will generate a fixture so that the valid record and PubMed ID can be reused in a later test.

\lstinputlisting[style=Python, firstline=22, lastline=29]{../tests/testing_chapter/test_textmining.py}
\lstinputlisting[style=Python, firstline=106, lastline=115]{../tests/testing_chapter/test_textmining.py}

Armed with these tests, we can then develop the code for \texttt{get\_record\_from\_PubMedID}:

\lstinputlisting[style=Python, firstline=48, lastline=65]{../src/bettercode/testing/textmining.py}

This passes the tests, so we can now move to writing some initial tests for \texttt{parse\_year\_from\_PubMed\_record}:

\lstinputlisting[style=Python, firstline=121, lastline=134]{../tests/testing_chapter/test_textmining.py}

And then we develop the implementation:

\begin{lstlisting}[style=Pythonshort]
def parse_year_from_PubMed_record(pubmed_record: dict) -> int:
    pubdate = pubmed_record.get("pubdate") 
    return int(pubdate.split()[0]) if pubdate else None
\end{lstlisting}

Now let's say that one had a specific set of PubMed IDs that they wanted to test the code against; for example, they might select IDs from papers published in various years across various journals.
To do this, we first create a list of tuples that include the information that we will need for the test; in this case it's the PubMed ID and the true year of publication.

\lstinputlisting[style=Python, firstline=140, lastline=148]{../tests/testing_chapter/test_textmining.py}

We then pass this into our test using the \texttt{@pytest.mark.parametrize} decorator on the test, which will feed in each of the values into the test successively:

\lstinputlisting[style=Python, firstline=151, lastline=156]{../tests/testing_chapter/test_textmining.py}

Note that we inserted a delay at the beginning of the test; this is necessary because the PubMed API will has a rate limit on requests, and running these tests without a delay to limit the request rate will result in intermittent test failures.
Looking at the results of running the test, we will see that each parametric value is run as a separate test:

\lstinputlisting[style=shell]{generated/testing/pytest_textmining_parametric.txt}

This test requires a live API, so it would fail in cases where one didn't have a proper network connection or if the API was down, and it would also be slow for a large number of tests.
It would be more efficient to mock the \texttt{get\_record\_from\_PubMedID} function to avoid dependency on the live API, but for our simple purposes it's fine to use the live API.

\section{Property-based testing}\index{testing!property-based}

Parameterized testing can be useful when we have specific values that we want to test, but sometimes we wish to test a large range of possible values drawn from some sort of distribution.
One approach to doing this is known as \textit{property-based testing}, and basically involves generating random values that match some specification and testing the code against those.

Property-based testing can be particularly useful for testing mathematical code, so we will develop another simple example to show how to use the \textit{hypothesis} module in Python to perform property-based testing.
Let's say that we have developed a function to perform linear regression, taking in two vectors (X and y variables) and returning a vector of length 2 (parameter estimates for slope and intercept).
Copilot generates some very terse code for us:

\begin{lstlisting}[style=Pythonshort]
def linear_regression(X, y):
    X = np.c_[np.ones(X.shape[0]), X]
    return np.linalg.inv(X.T @ X) @ X.T @ y
\end{lstlisting}

Asking Copilot to make the code more readable, we get this somewhat overly verbose version:

\lstinputlisting[style=Python, firstline=4, lastline=14]{../src/bettercode/testing/linear_regression.py}

The linear regression computation requires several things to be true of the input data in order to proceed without error:

\begin{itemize}
\item The data must not contain any infinite or \textit{NaN} values
\item The input data for X and y must each have at least two unique values
\item The X matrix must be full rank
\end{itemize}

In order to ensure that the input data are valid, we generate a function that validates the inputs, raising an exception if they are not, and we include this in our linear regression function:

\lstinputlisting[style=Python, firstline=17, lastline=36]{../src/bettercode/testing/linear_regression.py}

Now we can use the \texttt{hypothesis} module to throw a range of data at this function and see if it fails, using the following test:

\lstinputlisting[style=Python, firstline=8, lastline=22]{../tests/testing_chapter/_demos/test_property_based_smoke.py}

The \texttt{@given} decorator contains commands that will generate two arrays of the same size, which are then used as our X and y variables.
The main purpose of the test is to see whether the function successfully executes (i.e., a smoke test), but we include a minimal assertion to make sure that it returns a value that is not None.
We will turn off the validation in order to see what happens if the linear regression function is given invalid data.
Running this test, we see that the test fails, with the following output:

\lstinputlisting[style=shell]{generated/testing/pytest_property_based_smoke_fail.txt}

The test has identified a specific input that will cause the code to fail - namely, when the X variable is all zeros, which leads to an error when trying to invert the singular matrix.
We could get the test to pass by causing the function to return \texttt{None} when the matrix is not invertible, but this is not a great practice; we should announce problems loudly by raising an exception, rather than burying them quietly by returning \texttt{None}.

Now that we have seen how \textit{hypothesis} can identify errors, let's develop some tests for the code that we can use to make sure that it works properly.
We will first separately test the validator function, making sure that it can detect any of the potential problems that it should be able to detect:

\lstinputlisting[style=Python, firstline=10, lastline=30]{../tests/testing_chapter/test_property_based.py}

Note that this doesn't actually test whether our code actually gives the right answer, only that it runs without error and catches the appropriate problem cases, ensuring that any data passing the validator can run without error on the linear regression function.
When a reference implementation exists for a function (as it does in the case of linear regression), then we can compare our results to the results from the reference.
Here we will compare to the outputs from the from the linear regression function from the \textit{scipy} module.\index{SciPy}
Using this, we can check the randomly generated input to see whether it should raise an exception, and otherwise compare the results of our function to the \textit{scipy} function:

\lstinputlisting[style=Python, firstline=33, lastline=54]{../tests/testing_chapter/test_property_based.py}

This test passes, showing that our function closely matches the scipy reference implementation.
Note that we restricted the range of the values generated by the test to \texttt{[-1e6, 1e6]}; when the test values were allowed to vary across the full range of 64-bit floating point values (+/- 1.79e+308), we observed minute differences in the parameter estimates between the two functions that nonetheless exceeded the tolerance limits of \texttt{np.allclose()}.
We decided to restrict the test values to a range that is within the usual range of input data; if one were planning to work with very small or very large numbers, they would want to possibly test the input over a wider range and understand the nature and magnitude of differences in results between the methods.

\section{Automated testing and continuous integration}\index{continuous integration}

Once we have a set of tests for a project, it's important to integrate them into our workflow so that we ensure that new changes to the code don't break the existing code (known as a \textit{regression}).
One useful way to ensure that the tests are run regularly is to run them automatically every time changes are pushed to version control.\index{version control}
This is known as ``continuous integration'' (\textit{CI} for short), referring to the fact that it allows changes to be continuously integrated into the main code branch once they are confirmed to pass all of the tests.\index{Git!branching}
There are a number of different platforms one can use for CI; we will focus on \textit{Github Actions} since it is the most tightly integrated into the Github version control system.
Testing using CI also has a useful side effect: Since the CI system uses a virtual machine to run the tests, the use of CI for testing ensures that the code can run on a separate machine from the one where it was developed.\index{side effect}
Because setting up the CI system also requires understanding all of the dependencies that are required for the code to run, the CI setup provides a recipe to run the code on any other system.

\subsection{Using GitHub Actions}\index{GitHub}

GitHub Actions is a system that allows one to automatically execute workflows in response to events related to any GitHub repository.
At the time of writing, Github Actions are free and unlimited for public repositories when using the standard GitHub-based workflow runners.

When setting up an automated action using GitHub Actions, there are two primary decisions to specify:

\begin{itemize}
\item What is the workflow that I want to run?
\item What are the events that I want to trigger the workflow?
\end{itemize}

As an example, I will implement a workflow running to execute tests for a simple python package that was generated for this book project, called \textit{mdnewline} (github.com/poldrack/mdnewline).
I start by going to the GitHub Actions tab in the repository, and selecting the ``Python Package'' option, which creates a workflow that builds and tests a Python package.
Generating the workflow results in a file that contains a description of the workflow, located at \texttt{.github/workflows/python-package.yml}.
Looking more closely at the workflow file, we can see how it works.
The first section specifies the name of the workflow, and defines the events that will trigger the workflow:

\begin{lstlisting}[style=yaml]
name: Python package

on:
  push:
    branches: [ "main" ]
  pull_request:
    branches: [ "main" ]
\end{lstlisting}

This section specifies that the workflow will be run any time there is a push or a pull request to the main branch of the repo.\index{pull request}
The next section sets up workflow jobs:

\begin{lstlisting}[style=yaml]
jobs:
  build:
    name: python-build
    runs-on: ubuntu-latest
    strategy:
      matrix:
        python-version: ["3.12", "3.13"]
\end{lstlisting}

The \texttt{runs-on} argument tells the workflow runner which virtual machine to use to run the workflows; in this case, we will use the default which is the latest release of Ubuntu Linux.
We also tell it which versions of Python we want to test the code on, updating it to test on all current versions that are compatible with the package requirements (Python \textgreater = 3.12).

We then specify the actual steps in the workflow.
The first steps set up the workflow runner so that it can check out the repository, and set up the Python installation.
Since this project uses \texttt{uv} to manage packages, we will use the recommended setup code from the uv documentation for multiple Python versions:\index{uv}

\begin{lstlisting}[style=yaml]
steps:
     - uses: actions/checkout@v4

     - name: Install uv and set the python version
        uses: astral-sh/setup-uv@v6
        with:
          python-version: ${{ matrix.python-version }}
\end{lstlisting}

This installs uv with the appropriate Python version.
We can then install the project using \texttt{uv sync}, pip install the package within the \texttt{uv} environment, and run the tests:\index{pip}

\begin{lstlisting}[style=yaml]
- name: Install the project
        run: uv sync --locked --all -extras --dev

      - name: pip install the package
        run: uv pip install .

      - name: Run tests
        # For example, using *pytest*
        run: uv run pytest tests
\end{lstlisting}

When we commit and push this workflow file, it is automatically run by Github Actions.
% If we go to the Actions tab in the repository, we will see that the tests failed, and by looking at the logs we can see that the \texttt{uv} installation process failed:

% \includegraphics[width=0.7\linewidth]{files/github_actions_failu-45dfeb74c25ee1917e88387a5721a25c.png}

% After fixing the \texttt{uv} command, we now get the green light!:

% \includegraphics[width=0.7\linewidth]{files/github_actions_succe-81e11f1da04ae68f58bff1c62a826839.png}

% It's nice to advertise our testing to the world, which we can do by adding a status badge to our home page, using the markdown generated by Actions for us:

% \includegraphics[width=0.7\linewidth]{files/actions_status_badge-f7640e04ee103555b37619bdc6296c16.png}

\section{Optimizing the testing workflow}

As a project becomes larger and more complex, the tests will necessarily take longer to run - and whenever the tests take too long, it's likely that you won't run them as often.
This is particularly the case for data analysis tools, where testing the code on real data can take a very long time.
As an example, our lab develops a tool for fMRI data preprocessing called fMRIPrep \citep{Esteban:2019aa} that performs a large set of operations on functional MRI datasets.\index{fMRIPrep}
Without optimization, running the full test suite with real data would take roughly two hours, whereas with optimization we can complete it in 10-15 minutes.

There are a number of strategies to optimize one's testing workflows.

\subsection{Cherry-picking specific tests}

When developing a new function, it's usually sufficient to run only the tests that directly address that function rather than running the entire test suite.
If all of the tests for a specific function are located within a single test file, then one can simply call \textit{pytest} with that file.
It's also possible to run a specific test within a file by referring to a specific class or function using a double-colon marker:

\begin{lstlisting}[style=shell]
$ pytest tests/textmining/test_textmining.py::test_parse_year_from_PubMed_record

============================= test session starts ==============================
collected 1 item

tests/textmining/test_textmining.py .                                    [100%]
============================== 1 passed in 0.17s ===============================
\end{lstlisting}

This allows one to focus on the tests that are immediately relevant to a specific development task.

\subsection{Controlling test execution order}

Several of the strategies described below require the ability to execute tests in specific orders.
There are several ways one can use to do this.

One approach that \textit{will not} work reliably is to place the tests in a specific order in the test file.
While tests are often executed in the order that the functions appear in the file, this not guaranteed.

For precise ordering of particular tests, one can use the \textit{pytest-order} plugin, which allows the use of marks to specify test order\footnote{As discussed in the earlier section on technical debt, I don't think it's generally a good policy to rely upon packages that one randomly finds on Pypi or GitHub.
Before recommending \textit{python-order} as a possible solution, I looked at its GitHub page (github.com/pytest-dev/pytest-order), where I saw that it appears to be a well-maintained and currently active package, with recent commits and solid handling of issues.
Conversely, during the course of writing I came across a number of other packages that had been recommended on Stack Overflow to solve various problems, some of which had not seen commits in several years or had longstanding unaddressed issues.}.
We start with two tests that are out of order in the test code:

\begin{lstlisting}[style=Python]
import pytest

def test_second():
    """This is the second test."""
    assert True

def test_first():
    """This is the first test."""
    assert True
\end{lstlisting}

Running this, we see that these are run in the order they are defined in the test code:

\begin{lstlisting}[style=shell]
$ pytest -vv tests/ordering/test_order.py
=========================== test session starts ============================
collected 2 items

tests/ordering/test_order.py::test_second PASSED                     [ 50%]
tests/ordering/test_order.py::test_first PASSED                      [100%]
============================ 2 passed in 0.05s =============================
\end{lstlisting}

After installing the \texttt{pytest-order} plugin, we can now add marks that control the execution order of the tests:

\begin{lstlisting}[style=Python]
import pytest

@pytest.mark.order(2)
def test_second():
    """This is the second test."""
    assert True

@pytest.mark.order(1)
def test_first():
    """This is the first test."""
    assert True
\end{lstlisting}

Running this, we see that these are run in the order they are defined in the test code:

\begin{lstlisting}[style=shell]
$ pytest -vv tests/ordering/test_order.py
=========================== test session starts ============================
collected 2 items

tests/ordering/test_order.py::test_first PASSED                      [ 50%]
tests/ordering/test_order.py::test_second PASSED                     [100%]
============================ 2 passed in 0.05s =============================
\end{lstlisting}

This package also provides other useful ways to control execution order, such as relative markers like ``before'' or ``last''.

Note that while it might be tempting to use test ordering to ensure the proper executions of tests that rely upon one another, this is generally a bad practice.
Instead, one should generate fixtures for any dependencies that are shared by tests.

\subsection{Only run failed tests, or run them first}

Pytest has a set of features that allow one to focus testing on those tests that have most recently failed, which can be useful in the context of test-driven development.
The first command, \texttt{pytest --lf} (short for \texttt{--last-failed}) will check the cache of the most recent test run results, and only run those tests that failed in the most recent run.
It's important to realize that if one introduces a condition that would cause a previously passed test to fail, it may not be caught, such that one could run a string of tests using the \texttt{--lf} flag, get to the point where the tests all seem to pass, but then find that some tests that had passed earlier will fail due to conditions introduced when fixing the later tests.

A related command, \texttt{pytest --ff} (for \texttt{--failed-first}) will run the most recent failed tests before running the rest of the tasks.
This command is probably better for general use than the \texttt{--lf} flag, since it ensures that all tests are run, but prioritizes the running of the failing tests first.
This might be an option that one would consider using by default, which can be achieved by adding the option to one of the relevant configuration files.
In this case the relevant file is \texttt{pyproject.toml}, where we could add:

\begin{lstlisting}[style=toml]
[tool.pytest.ini_options]
addopts = " --ff"
\end{lstlisting}

\subsection{Separate unit tests from longer-running tests}

If we are writing short functions (which we are, right?), then our unit tests should usually run fairly quickly, whereas our integration tests (which integrate across a potentially large number of modules) may take much longer.
One thing we can do is to run only unit tests while we are developing, while having the integration tests run less frequently.
For example, we could have the unit tests run any time we commit our code (for example, by running them as pre-commit items), while having the integration tests run automatically on our CI system when the code is actually pushed to the GitHub repository.\index{pre-commit}

One way to accomplish this is to use \textit{pytest} markers (as we saw earlier) to mark groups of tests.
We first need to define our own custom markers in our \texttt{pyproject.toml} file:

\begin{lstlisting}[style=toml]
[tool.pytest.ini_options]
markers = [
    "unit: mark a test as a unit test.",
    "integration: mark a test as an integration test."
]
\end{lstlisting}

We can then generate tests using these markers:

\begin{lstlisting}[style=Python]
import pytest

@pytest.mark.unit
def test_unit1():
    """This is a unit test."""
    assert True

@pytest.mark.unit
def test_unit2():
    """This is a unit test."""
    assert True

@pytest.mark.integration
def test_integration():
    """This is an integration test."""
    assert True
\end{lstlisting}

Running these using the standard pytest command, both tests are executed:

\begin{lstlisting}[style=shell]
$ pytest -vv tests/markers/test_markers.py
=========================== test session starts ============================
collected 3 items

tests/markers/test_markers.py::test_unit1 PASSED                     [ 33%]
tests/markers/test_markers.py::test_unit2 PASSED                     [ 66%]
tests/markers/test_markers.py::test_integration PASSED               [100%]
============================ 3 passed in 0.05s =============================
\end{lstlisting}

However, we can also specify that only the unit tests should be executed, using the \texttt{-m} flag to deselect all other flags:

\begin{lstlisting}[style=shell]
$ pytest -vv -m unit tests/markers/test_markers.py
=========================== test session starts ============================
collected 3 items / 1 deselected / 2 selected

tests/markers/test_markers.py::test_unit1 PASSED                     [ 50%]
tests/markers/test_markers.py::test_unit2 PASSED                     [100%]
===================== 2 passed, 1 deselected in 0.05s ======================
\end{lstlisting}

We could also run all tests that to not match a particular marker:

\begin{lstlisting}[style=shell]
$ pytest -vv -m "not unit" tests/markers/test_markers.py
=========================== test session starts ============================
collected 3 items / 2 deselected / 1 selected

tests/markers/test_markers.py::test_integration PASSED               [100%]
===================== 1 passed, 2 deselected in 0.05s ======================
\end{lstlisting}

\subsection{Longest tests last}

Another simple strategy that can help optimize the testing workflow is to run tests in order to the time for completion.
This ensures that one doesn't end up waiting a long time for long-running tests to complete, only to find that a quick test fails.

We can use pytest's \texttt{--durations} flag to get the durations for each test.
Here is an example with three tests with different durations to completion:

\begin{lstlisting}[style=Python]
import pytest
from time import sleep

def test_duration_3():
    sleep(3)
    assert True

def test_duration_5():
    sleep(5)
    assert True

def test_duration_1():
    sleep(1)
    assert True
\end{lstlisting}

We can call pytest with \texttt{--durations=0} which will return the duration for all of the tests:

\begin{lstlisting}[style=shell]
$ pytest --durations=0 tests/ordering/test_duration.py
============================= test session starts ==============================
collected 3 items

tests/ordering/test_duration.py ...                                      [100%]

============================== slowest durations ===============================
5.00s call     tests/ordering/test_duration.py::test_duration_5
3.01s call     tests/ordering/test_duration.py::test_duration_3
1.00s call     tests/ordering/test_duration.py::test_duration_1
============================== 3 passed in 9.10s ===============================
\end{lstlisting}

Using the results of this, we could impose ordering on test execution so that the slowest ones are executed last, or add a \texttt{slow} marker that we could exclude from our regular testing.

There are plugins for \textit{pytest} that can measure the duration of each test and order tests accordingly, but they don't seem to be very well maintained so I am avoiding them.
However, with agentic AI coding tools we can have our AI system make the appropriate call to \textit{pytest} to obtain the durations and then add the appropriate ordering markers.
I tried this using the agentic chat window in VSCode (with Claude Sonnet 4).\index{VSCode}\index{Claude}
On the first pass, Claude immediately noticed that it could infer the completion times directly from the \texttt{sleep()} commands in the code.
I asked it not to do this, and instead to use the outputs from \texttt{pytest --durations} but it had trouble waiting for the tests to finish, and ended up continuing to try to things that didn't work.
However, Claude Code was able to successfully solve the problem; the first line marked with ``\textgreater '' was my prompt, and the remainder is Claude Code's output:

\begin{lstlisting}[style=transcript]
> the file tests/ordering/test_duration.py contains several
  tests.  Use pytest --duration to determine the durations for
  each of the tests in this file. DO NOT try to infer the 
  durations from the code itself. The tests will take about
  fifteen seconds to complete, so be sure to wait for them. 
  Then, add marks to each test using pytest-order to specify 
  that the fastest tests are executed first.

⏺ I'll run pytest with the --durations flag to measure the 
actual test durations, then add pytest-order marks to run the
fastest tests first.

⏺ Update Todos
  ⎿  ☐ Run pytest with --durations flag on test_duration.py
     ☐ Analyze test duration results
     ☐ Add pytest-order marks to tests based on duration

⏺ Bash(pytest tests/ordering/test_duration.py --durations=0 -v)
  ⎿ ============================= test session starts =====================
    =========
    platform darwin -- Python 3.12.0, pytest-8.4.1, pluggy-1.5.0 -- /Users/
    … +27 lines (ctrl+r to expand)

⏺ Based on the test durations:
  - test_duration_1: 1.00s (fastest)
  - test_duration_3: 3.00s
  - test_duration_5: 5.01s (slowest)

⏺ Now I'll add pytest-order marks to run tests from fastest to slowest:

⏺ Update(tests/ordering/test_duration.py)
  ⎿  Updated tests/ordering/test_duration.py with 3 additions
        2    from time import sleep
        3
        4
        5 +  @pytest.mark.order(2)
        6    def test_duration_3():
        7        sleep(3)
        8        assert True
        9
       10 +  @pytest.mark.order(3)
       11    def test_duration_5():
       12        sleep(5)
       13        assert True
       14
       15 +  @pytest.mark.order(1)
       16    def test_duration_1():
       17        sleep(1)
       18        assert True

⏺ Let me verify the ordering works correctly:

⏺ Bash(pytest tests/ordering/test_duration.py -v)
  ⎿ ============================= test session starts =====================
    =========
    platform darwin-- Python 3.12.0, pytest-8.4.1, pluggy-1.5.0 -- /Users/
    … +20 lines (ctrl+r to expand)

⏺ Perfect! The tests now run in order from fastest to slowest:
  1. test_duration_1 (1 second)
  2. test_duration_3 (3 seconds)
  3. test_duration_5 (5 seconds)
\end{lstlisting}

This is a great example of how we can increasingly rely upon AI coding models to solve problems that would once have required specialized tools.

\subsection{Using minimal mock datasets}

For code that performs data processing operations, the processing of full size datasets can often take a very long time.
One strategy in these cases is to generate minimal mock datasets that can exercise the functions without taking the full amount of time that a real dataset would.
In the tests for \textit{fMRIPrep}, we use fMRI datasets that have been reduced in length, and structural MRI datasets that have been downsampled to reduce their spatial resolution.
The specific way to reduce the dataset will depend on the particular processes being run.
For example, downsampling the data too much for MRI preprocessing would likely cause some operations to fail, so one needs to have a good intuition for the data requirements of the relevant code.

\subsection{Adding minimal processing modes for integration tests}

When the goal is to the test the integration of components rather than the function of each component, one way to minimize testing time is to provide configuration features that minimize execution time for the component.
For example, in \textit{fMRIPrep} there are a number of steps that involve optimization processes that take time to converge.
However, there the package has a ``sloppy mode'' configuration flag that one can turn on for testing, which provides a more lenient threshold for convergence of those operations, allowing them to finish faster.
Again, knowing where one can cut corners requires a good understanding of the specific requirements of the processing operations.

\subsection{Parallelizing testing}

If we have written good tests, they should be able to run independently, and thus their execution should be parallelizable, assuming that we are using a system with multiple CPU cores.
If we are using \textit{pytest} as our testing framework, then we can use the \textit{pytest-xdist} extension to enable the parallel execution of tests in pytest.
For example, let's set up a parameterized test that includes a \texttt{time.sleep()} command so that execution will take a significant amount of time.

\begin{lstlisting}[style=Python]
import pytest
import time

@pytest.mark.parametrize("x", range(10))
def test_parallel(x):
    assert x in range(10), f"Value {x} is not in the expected list."
    time.sleep(1) # wait for one second
\end{lstlisting}

If we run this using the standard \texttt{pytest} command, we should see that it takes about ten seconds, given that there are ten tests:

\begin{lstlisting}[style=shell]
$ pytest tests/parallel/tests_parallel.py
======================== test session starts =========================
collected 10 items

tests/parallel/tests_parallel.py ..........                     [100%]
======================== 10 passed in 10.18s =========================
\end{lstlisting}

If we have installed \texttt{pytest-xdist} then we can add the \texttt{-n auto} flag which will automatically detect how manu CPU cores we have available and run the tests in parallel across those cores:

\begin{lstlisting}[style=shell]
$ pytest tests/parallel/tests_parallel.py -n auto
========================= test session starts ==========================
16 workers [10 items]
..........                                                        [100%]
========================== 10 passed in 1.97s ==========================
\end{lstlisting}

You can see that it detected the 16 cores in my laptop and ran the 10 tests in parallel, greatly reducing the testing time.